\documentclass[10pt,oneside]{article}

%%%%%%%%%%%%%
\setlength{\textheight}{8.75in} %Letter is 11in, less 2 for margins, less 0.25 for footer
\setlength{\oddsidemargin}{0.0in} %gets +1inc
\setlength{\evensidemargin}{0.0in} %gets +1inch
\setlength{\textwidth}{6.50in} %Letter is 8.5, less 2 inches for margins
\setlength{\topmargin}{0.5in}
\setlength{\headheight}{0in}
\setlength{\headsep}{0in}
\setlength{\parindent}{0.25in}
%%%%%%%%%%%% 

\usepackage[numbers]{natbib}
\usepackage{graphicx}
\usepackage[colorlinks=true,citecolor=black,urlcolor=blue]{hyperref}

\title{%Hack to get the logo on the PDF front page:
\vspace{-1.5in}
\includegraphics[width=0.5\textwidth]{biopython.jpg}\\
%Hack to get some white space using a blank line:
~\\Biopython Project Update 2014}
\author{
    \underline{Wibowo Arindrarto}\footnote{Sequencing Analysis Support Core, Leiden University Medical Center, Leiden, NL. Email: \href{mailto:bow@bow.web.id}{bow@bow.web.id}},
    Peter Cock\footnote{Information and Computational Sciences, James Hutton Institute (formerly SCRI), Invergowrie, Dundee, UK},
    Eric Talevich\footnote{Department of Dermatology, University of California San Francisco, San Francisco, CA, USA},
    Michiel de Hoon\footnote{Division of Genomic Technologies, RIKEN Center for Life Science Technologies, Yokohama, JP},
    Tiago Antao \footnote{Vector Biology Department, Liverpool School of Tropical Medicine, Pembroke Place, UK},
    \\
    Jo\~{a}o Rodrigues\footnote{Bijvoet Center for Biomolecular Research, Utrecht University, Utrecht, NL},
    \textit{et al.}}
\date{15\textsuperscript{th} Bioinformatics Open Source Conference (BOSC) 2014, Boston, MA, USA}

\begin{document}
%\pagestyle{empty}
%\thispagestyle{empty}
\maketitle
%\pagestyle{empty}
\thispagestyle{empty}

\vspace{-0.2in}
\noindent
Website: \url{http://biopython.org} \\
Repository: \url{https://github.com/biopython/biopython} \\
License: Biopython License Agreement (MIT style, see \url{http://www.biopython.org/DIST/LICENSE}) \\

We present the latest updates from the Biopython project, a long-running,
distributed collaboration producing a freely available Python library for
biological computation \citep{AppNote}. Biopython is supported by the Open
Bioinformatics Foundation (OBF).

There have been three releases since BOSC 2013: version 1.62, 1.63, and 1.64; all
involving contributions made by new and returning developers. New features
in version 1.62 include parsing support for NeXML and CDAO in the Bio.Phylo
module, parsing support for GAF, GPA, and GPI formats from UniProt-GOA
in the Bio.UniPort module, and BioSQL support for Jython, among others.
In version 1.63, we added  support for the population genetic tool
fastsimcoal, a wrapper for samtools, and other significant enhancements to
existing modules. Version 1.64 saw the addition of the Bio.CodonAlign module
and enhancements to the Bio.Phylo module, contributed by our Google Summer
of Code (GSoC) 2013 students. The upcoming version 1.65 is now under development.
Moreover, since BOSC 2013 we have successfully supported Python 2, Python 3,
PyPy, and Jython 2.7 with a single codebase. This change is also reflected in
our Tutorial \& Cookbook, which uses code compatible with all Python versions.

In addition to local installation on various operating systems, Biopython
is now available at \url{http://toolshed.g2.bx.psu.edu/view/biopython/}
in the Galaxy Tool Shed \citep{ToolShed} as a package dependency. Galaxy tools
requiring Biopython can now specify this dependency explicitly and choose from
Biopython version 1.63 onwards.

We participated in GSoC 2013 under the umbrella of the National Evolutionary
Synthesis Center (NESCent). Two students were selected: Yanbo Ye, working on
enhancements for Bio.Phylo, and Zheng Ruan, working to add codon alignment and
analysis support. Both students have finished their project successfully, with
Yanbo's code integrated into the existing Bio.Phylo module and Zheng's code
submitted as a pull request under review. We continue to participate in GSoC 2014
with the selection of Evan Parker who will be working on adding lazy-parsing
support to Bio.SeqIO.


\begin{thebibliography}{}

\bibitem[Cock {\it et al}., 2009]{AppNote}Cock, P.J.A., Antao, T., Chang, J.T., Chapman, B.A., Cox, C.J., Dalke, A., Friedberg, I., Hamelryck, T., Kauff, F., Wilczynski, B., de Hoon, M.J. (2009) Biopython: freely available Python tools for computational molecular biology and bioinformatics. {\it Bioinformatics} {\bf 25}(11) 1422-3. \href{http://dx.doi.org/10.1093/bioinformatics/btp163}{doi:10.1093/bioinformatics/btp163}
\bibitem[Blankenberg {\it et al}., 2014]{ToolShed}Blankenberg, D., Von Kuster, G., Bouvier, E., Baker, D., Afgan, E., Stoler N., The Galaxy Team, Taylor, J., Nekrutenko, A. (2014) Dissemination of scientific software with Galaxy ToolShed. {\it Genome Biology} {\bf 15} 403. \href{http://dx.doi.org/10.1186/gb4161}{doi:10.1186/gb4161}

\end{thebibliography}

\end{document}
